\documentclass[11pt,a4paper]{article}
\usepackage[T1]{fontenc}
\usepackage{lmodern,amsmath,amssymb,geometry}
\usepackage[hidelinks]{hyperref}
\geometry{margin=2.3cm}
\setlength{\parindent}{0pt}
\setlength{\parskip}{0.5em}
\title{BM4Midpoint: Arithmetic Diagonal Projection}
\author{GC2D numerical methods}
\date{}
\setlength{\emergencystretch}{3em}
\begin{document}
\maketitle

\section*{Explicit palindromic operations}
All ABBA and BM4 presets inherit \texttt{PalindromicMethod}. For signed
$d=c h$, the elementary operations are
\[
 A_d(u,v,t)=(u+d f(t,v),v,t),\qquad
 B_d(u,v,t)=(u,v+d f(t,u),t),\qquad
 T_d(u,v,t)=(u,v,t+d).
\]
Only $T$ advances time. A coefficient $c=1/2$ means duration $h/2$;
there is no additional half inside $A$ or $B$. The optional exact $C$ flow
rotates the copy difference and requires an independent coupling object.
Both BM4 presets omit $C$.

In chronological order, standard ABBA2 is
$A_{h/2},B_{h/2},T_h,B_{h/2},A_{h/2}$.
The distinct symmetric variant
$T_{h/2},A_{h/2},B_{h/2},B_{h/2},A_{h/2},T_{h/2}$
evaluates every spatial shear at the midpoint time. Both exhibit second
order on the tested nonautonomous rotation. Arbitrary palindromes do not
inherit fourth or sixth order, and an arithmetic projection need not be
reversible or symplectic.

\texttt{Composition} denotes an explicit tuple of operations.
Consecutive time translations combine as $T_a T_b=T_{a+b}$.
The generic \texttt{PalindromicImplicit} and \texttt{PalindromicMidpoint}
methods select one exterior Hairer projection or a final arithmetic mean.
The signed spatial flows and passive momentum normalization are retained;
observers reconstruct their groups from elementary traces without replay. NumPy and JAX execute the same
explicit operations and differentiate their retained shear inputs.
See \path{docs/models/extended/simulation/palindromic-operations.md}.

\section*{Current uncoupled configuration}
Both \texttt{BM4Implicit} and \texttt{BM4Midpoint} omit the harmonic
coupling flow and no longer accept a coupling-frequency constructor argument.
Their twelve coefficients, signed clock and projection placement are unchanged.
The independent operation \texttt{HarmonicCoupling} retains C for future
compositions. Existing observer/archive rates are fixed at zero for new runs;
nonzero historical calculations require their original frozen implementation.

\section{Complete-cycle definition}
For $\dot z=f(t,z)$, embed the accepted state as $E z=(z,z)$ and define
$P(u,v)=(u+v)/2$. One step is
\[
 z_{n+1}=P\Psi_h(Ez_n),
\]
where $\Psi_h$ is the same unprojected twelve-stage BM4 composition used by
BM4Implicit. No projection is applied between stages and no multiplier is
solved. ``Midpoint'' refers to the final arithmetic mean.

The coefficient vector is $(b_1,\ldots,b_6,b_6,\ldots,b_1)$, where
\begin{align*}
b_1&=0.0792036964311957,& b_2&=0.1303114101821663,\\
b_3&=0.2228614958676077,& b_4&=-0.3667132690474257,\\
b_5&=0.3246481886897062,& b_6&=0.1096884778767498.
\end{align*}
The sum is one to floating-point accuracy. Starting with the adjoint, stages
alternate adjoint and direct maps with signed duration $d_j=b_jh$.
If $s_j=t_n+\sum_{i<j}d_i$, an adjoint stage uses $s_j$ and a direct stage uses
$s_j+d_j$. The negative coefficient also moves the stage clock backwards.

At the chosen evaluation time, let
\[
A_d(u,v)=(u+d f(t,v),v),\qquad
B_d(u,v)=(u,v+d f(t,u)).
\]
The spatial direct map is $A_d\circ B_d$ and its adjoint is
$B_d\circ A_d$. Neither includes harmonic coupling. The time translations
remain in the complete direct/adjoint stage as specified above.
Each stage evaluates the vector field twice.
Thus one complete step uses 24 evaluations and one arithmetic projection.

\section{State and energy conventions}
Twelve direct/adjoint stages followed by one spatial arithmetic mean. Energy tracking is optional and passive.
The two shear contributions in each stage update a summed momentum. Dividing this accumulator by two gives the physical kappa normalization.
The common state conventions are specified below.

\section{Order and geometric limitations}
On the invariant diagonal of the exact extended flow, the two copies follow
the original dynamics and the harmonic coupling vanishes. For a smooth field,
the fourth-order unprojected composition therefore satisfies
\[
 \Psi_h(Ez)=E\varphi_h(z)+O(h^5).
\]
Since $P E=I$ and $P$ is linear, the projected method has local error
$O(h^5)$ and global order four under the usual stability assumptions.
The passive energy coordinates do not change the physical map or its order.

Arithmetic projection discards the copy difference. Consequently, symmetry
or symplecticity of a base composition alone does not establish exact
reversibility or symplecticity of the accepted physical map. This method
makes neither claim. The recorded infinity norm $\|u_f-v_f\|_\infty$ measures
the discarded separation before projection.

\section{Executable checks}
The test suite compares the physical map with an independent signed-stage
replay and checks fourth-order convergence against an exact rotation and a
tightly integrated nonautonomous reference with energy tracking both disabled and enabled.
It also checks momentum normalization against an augmented reference ODE,
particle batching, persistent observer callbacks, and independence from
off-grid output sampling. Nonlinear iterations are absent by construction.
\section{Physical accepted states and passive energy tracking}
For $N$ particles with $d$ physical components, every accepted snapshot contains
$Z\in\mathbb{R}^{dN}$ and one scalar time $t$. An optional energy observer
owns separate physical momentum $\kappa\in\mathbb{R}^{N}$.
Thus GC accepted vectors have $2N$
components and FC vectors have $4N$, whether energy is tracked or not.
A doubled method prepares temporary $(u,v)\in\mathbb{R}^{2dN}$ copies,
allows them to separate, and projects to one physical result before acceptance.
The representation stores only the physical size and never owns a clock,
momentum, physical coordinate declaration, dynamics or solver.

The normalized passive balance is
\[
\dot\kappa=-\partial_tH(t,Z),\qquad \kappa(t_0)=0,\qquad
\Delta K(t)=H(t,Z(t))+\kappa(t)-H(t_0,Z_0).
\]
GC split-map contributions are summed in temporary work and divided by two
to obtain physical momentum. Tracking leaves physical stages, nonlinear
stopping scales and adaptive acceptance unchanged. An autonomous Hamiltonian
supplies an explicit zero momentum derivative.

Hairer's constraint is solely $u-v=0$: the reduced multiplier has $2N$ entries
and ABBA's simultaneous spatial solve has $6N$ unknowns. Neither $t$ nor
$\kappa$ is projected. The former eight-component projection is retired.
Dynamics declare component order and position/velocity roles through immutable
structural data; initial conditions and results retain those data independently.
Execution injects $f(t,Z)$ and optional physical derivatives into a fresh method
run. Methods own physical stages and projection equations, without retaining the
dynamics object. Nonlinear tolerances, iteration limits, differentiation and
iteration callbacks belong to the configured algebraic solver. Generic doubled
copy operations live in \path{src/formulations/state.py}; classical, HBVM and
adaptive methods need no state formulation.

Methods optionally return retained stage quadrature or a dense interpolant.
The energy observer integrates momentum independently, using Hamiltonian
functions supplied by execution. Controllers keep only physical accepted states;
compiled JAX loops carry the observer accumulator separately.
Observers receive physical states. Energy histories in \texttt{States.diagnostics}
retain their existing labels; \texttt{extended\_time} has shape
$(N,n_{\rm samples})$ and is broadcast from scalar sample times.
Fixed shadow samples leave accepted states and counters unchanged. Adaptive
quadrature uses accepted dense output and physical-only error control.




\section{Shared extended-space implementation}
Public method names remain unchanged; both BM4 constructors now omit harmonic
coupling. Their code
lives in \path{src/methods/extended/}. Public exports from \texttt{simulation}
remain available; former ABBA and BM4 import adapters have been removed.
The shared execution is
\texttt{advance}, \texttt{solve\_projection}, \texttt{compose}, then
explicit $A/B/T/C$ operations. Arithmetic midpoint variants
replace the nonlinear projection by one final average.

The composition retains signed shear inputs for exact particle-block tangents,
passive energy quadrature, and existing observation records. Newton and Broyden
share the spatial projection equations; the energy observer accumulates only the
accepted trace, with the existing factor $1/2$ for physical $\kappa$.
Adjoint stages evaluate at the start of their signed substep, direct stages at
its end. The coefficients remain unchanged. ABBA2Implicit, ABBA4Implicit
and BM4Implicit each use one outer projection around their complete tuples,
and directly inherit \texttt{ProjectedPalindromicMethod}. ABBA6 is removed;
its saved results retain their historical labels. The internal family name
does not enable the retired full time/momentum projection.

The implementation architecture and current diagrams are documented in
\path{docs/models/extended/simulation/extended-simulation-architecture.md}.


\section{Optional JAX execution}

An optional JAX driver executes the same numerical map on one selected CPU
or GPU, with particle operations batched inside a compiled sequential time
loop. It reuses the canonical coefficients and prepared potential splines.
Physical-state tolerance scaling, signed stage times, passive energy weights
and independent shortened output maps retain their definitions above.
Applicable Newton systems use independent particle blocks; an explicitly
selected Broyden iteration retains its global secant matrix. Nonlinear
iteration limits and failure checks remain active on the device.

Both execution routes use double precision and return the same physical
solution layout. Floating-point reassociation can produce round-off
differences without changing the designed order or geometric properties.
First-call compilation is separate from warmed integration timings, which
include transfers and output construction. The sibling runtime document and
shared execution contract specify supported callbacks and exact device scope.

\end{document}
